\section{Background and terminology}
\label{sec:background}

\begin{inbrief}
This section recalls the standard notions the paper uses: task-chains
and their end-to-end metrics, fixed-priority scheduling, LET, TSN and
Network Calculus. \Cref{tab:glossary} lists the terms.
\end{inbrief}

\subsection{Task-chains and their end-to-end metrics}

A \emph{task} is released periodically and runs for at most its
worst-case execution time (WCET). A \emph{task-chain} is a sequence of
tasks in which each task reads the output of the previous one. Automotive
task-chains are multi-rate: the tasks of a chain may have different
periods. Consecutive tasks with the same period form an \emph{activation
pattern}. They are released together and run in order. Between two
patterns, the tasks communicate through a shared variable (a register):
the consumer reads the latest value at its own release
\cite{kramer2015waters,hamann2017communication}.

Two end-to-end metrics are standard \cite{feiertag2008compositional,
becker2017end,durr2019end}. The \emph{maximum reaction time} is the
longest time from an external event (a change seen by the sensor) to the
end of the first job of the last task whose output accounts for it. The
\emph{maximum data age} is the longest time a sensor value is used, from
its sampling to the end of the last job of the last task that uses it.
Günzel et al.~\cite{gunzel2023equivalence} proved that the two are equal
for cause-effect chains of periodic tasks. In this paper, the end-to-end deadline of a
task-chain bounds its maximum reaction time, and the bound therefore also
holds for the data age. The \emph{end-to-end response time} of holistic
analysis \cite{tindell1994holistic} is a third quantity: the sum of the
response times and network delays along the chain, without the sampling
and copy delays. It is one part of the reaction time, not a time of the
system.

\subsection{Fixed-priority scheduling on partitioned multicores}

Each core runs its tasks with preemptive fixed priorities, and a task
never migrates (partitioned scheduling \cite{davis2011survey}). The
worst-case response time (WCRT) of a task is the longest time from its
release to its completion; it is found by the classical response-time
analysis \cite{joseph1986finding,audsley1993applying}. A core is
\emph{schedulable} if every task meets its implicit deadline, its period.
Audsley's algorithm \cite{audsley2001priority} assigns the priorities
from the lowest level up; for schedulability tests that depend only on
the set of higher-priority tasks, such as response-time analysis, it finds
a schedulable order whenever one exists \cite{davis2016review}. The \emph{critical scaling factor} of a core is the largest
factor by which all WCETs can be multiplied with the core still
schedulable \cite{lehoczky1989rate}; it measures the timing margin.
Allocating tasks to cores under a utilization cap is a bin-packing
problem; the classical heuristics take the tasks by decreasing
utilization and put each one on the first (first fit), fullest (best fit)
or emptiest (worst fit) core that can take it \cite{johnson1974fast,
davis2011survey}.

\subsection{Logical execution time}

With the LET paradigm \cite{henzinger2003giotto,kirsch2012let}, a task
reads its inputs at its release and publishes its outputs at its next
release (or deadline), whatever its actual execution. Between cores, the
copies are made by a communication task at fixed instants
\cite{hamann2017communication,biondi2018achieving}. The timing of the data
flow then no longer depends on execution times, at the price of sampling
delays and of the copy task. With \emph{implicit communication}, the
AUTOSAR alternative, a task works on local copies of its data, which are
written back when it finishes \cite{hamann2017communication}: the delays
are shorter, but they depend on response times. System-level LET extends
LET across ECUs, with synchronized clocks \cite{gemlau2021system}.

\subsection{Time-Sensitive Networking}

TSN is a set of IEEE~802.1 standards for bounded-latency Ethernet
\cite{ieee8021q2022}. A switch has one output queue per traffic class at
each port. The paper uses four mechanisms.
\begin{itemize}
\item \emph{Credit-based shaper} (CBS, 802.1Qav): a class may send only
while its credit is non-negative; the credit grows at the \emph{idle
slope} while frames wait and drains while the class sends. It reserves a
share of the link for the class and spreads its frames out.
\item \emph{Time-aware shaper} (TAS, 802.1Qbv): a cyclic gate schedule
opens the queue of each class in its own time window.
\item \emph{Asynchronous traffic shaping} (ATS, 802.1Qcr
\cite{ieee8021qcr2020,specht2016urgency}): at every switch, a
\emph{regulator} holds each frame until the stream again conforms to its
token-bucket description, and the classes are then served in strict
priority. Bursts therefore do not grow from hop to hop.
\item \emph{Frame replication and elimination for reliability} (FRER,
802.1CB \cite{ieee8021cb2017}): a stream is sent as two \emph{member
streams} over disjoint paths; the receiving side keeps the first copy of
each frame and discards the other, so a single link or switch failure
loses no frame.
\end{itemize}
Frame preemption (802.1Qbu) lets a time-critical frame interrupt a
best-effort frame \cite{thiele2016formal}.

\subsection{Network Calculus}

Network Calculus bounds worst-case delays with curves
\cite{cruz1991calculus,leboudec2001network}. An \emph{arrival curve}
$\alpha$ bounds the traffic of a flow: in any window of length $t$, the
flow sends at most $\alpha(t)$ bits. The token bucket $\alpha(t)=b+rt$
(burst $b$, rate $r$) is the usual one. A \emph{service curve} $\beta$
guarantees the service a port gives to a class; the rate-latency curve
$\beta(t)=R\,[t-L]^+$ (rate $R$ after a latency $L$) is the usual one.
The delay of a flow at the port is then at most the horizontal distance
between $\alpha$ and $\beta$, here $L+b/R$; and the flow leaves the port
with a larger burst, $b+rd$ for a delay bound $d$. Along a route, bursts
therefore grow from port to port, unless regulators reshape them.

\begin{table}[t]
\centering\footnotesize
\caption{Terms used in the paper.}
\label{tab:glossary}
\setlength{\tabcolsep}{3pt}
\rowcolors{2}{brandlight}{white}
\begin{tabularx}{\columnwidth}{@{}lX@{}}
\toprule
\rowcolor{white}
term & meaning \\
\midrule
task-chain & sequence of tasks, sensor to actuator (cause-effect chain) \\
activation pattern & maximal run of consecutive tasks of a chain with one period \\
sub-chain & run of consecutive tasks of a chain, the unit of allocation \\
segment & maximal run of consecutive sub-chains of a chain on one core \\
pinned & bound to one ECU (sensor and actuator drivers) \\
central computer & the ECU that runs most functions \\
zone controller (ZCU) & ECU of one zone of the car, wired to its devices \\
host & ECU that free sub-chains go to first: the central computers, else the ECUs with the fastest cores \\
inter-ECU communication & data sent between two segments on two ECUs \\
LET task & copy task of a core, run at its LET instants \\
implicit / strict LET & output copied after its producer ends / at the end of its period \\
utilization cap $\varphi$ & largest load allowed on a core \\
maximum reaction time & event to the first actuation that accounts for it \\
end-to-end response time & sum of the response times and network delays of a chain (holistic) \\
critical scaling factor & largest WCET factor that keeps the cores schedulable \\
design point & one run of the method for one application and one set of parameters \\
SOP & start of production: the hardware is frozen \\
\bottomrule
\end{tabularx}
\end{table}
